\begin{figure}[htbp]
\centering
\begin{tikzpicture}[x=1cm,y=1cm,>=stealth]
\fill[AnalysisBlue!9] (-3,0)--(0,3)--(3,0)--cycle;
\draw[->] (-3.8,0)--(3.8,0) node[right]{\(x\)};
\draw[->] (-3.6,-.2)--(-3.6,3.5) node[above]{\(t\)};
\draw[thick,AnalysisBlue] (-3,0)--(0,3)--(3,0);
\draw[dashed] (-1.8,1.2)--(1.8,1.2);
\fill (0,3) circle(1.7pt) node[above right]{\((t_0,x_0)\)};
\node[below] at (-3,0){\(x_0-ct_0\)};
\node[below] at (3,0){\(x_0+ct_0\)};
\node[below] at (0,0){\(x_0\)};
\node at (0,.6){\(K_{t_0,x_0}\)};
\node[above] at (0,1.2){时刻 \(s\) 的截面};
\end{tikzpicture}
\caption[一维波动方程的依赖三角形]{一维波动方程的依赖三角形。截面半径为 \(c(t_0-s)\)，两侧边界随时间以速度 \(c\) 向观察点收缩；图中横纵比例仅为示意。}
\label{b2:fig:09-cone}
\end{figure}
